\section{The frame quotient and compact supports}
\label{sec:support}

This section converts the actual Frobenius transpose of
Lemma~\ref{lem:cartier-transfer} into a calculation of compact supports on a
space of independent bundle extensions. The calculation proves finiteness
of the stable transfer image in Corollary~\ref{support:finite-stable-image},
which Section~\ref{sec:compact} promotes to finiteness for compact coefficients.

Throughout this section, $1\leq t<n<p-1$ and $m=n-t$.
We use the rings, extendable frame levels, and transfer of
Section~\ref{sec:rings}.  In particular, $x=x_t$, the first labelled
ring is $C_1$, and the joint frame group is
\[
 U_0=P_t\times\GL_m(\Zp).
\]
Let $C$ be a complete algebraically closed extension of the unramified
ground field, let $E=C^\flat$, and choose compatible roots of $p$ in $C$.
Their tilt is denoted by $\varpi$, so $\varpi^\sharp=p$.  Set
\[
 R=(\cO_E/\varpi)^a,
 \qquad
 \mathscr A=(\cO^{\flat,+}/\varpi)^a.
\]
The superscript $a$ means almostification with respect to the maximal
ideal of $\cO_E$.  On a diamond over the chosen untilt, $\mathscr A$
identifies with $(\cO^+/p)^a$; this identification is compatible with
pullback.  We retain Tate twists in intermediate formulas.  A choice of
compatible roots of unity over $C$ identifies $R(1)$ with $R$ whenever
only the frame-group actions are under consideration.

\subsection{From Tate coordinates to extensions of bundles}

Write $V_a=\cO(1/a)$ for the stable bundle of rank $a$ and degree one
on the Fargues--Fontaine curve.  These conventions agree with the
unramified-pushforward definition of the standard bundles
\cite[Definition~8.2.2 and Theorem~8.2.10]{FarguesFontaine}.
Define the relative extension sheaf
\[
 N_t=\mathcal{E}xt^1(V_t,\cO),
 \qquad Z=(N_t^m)_{\mathrm{ind}}.
\]
Here $\mathcal{E}xt^1$ is sheafified on perfectoid test spaces.  A tuple
$(e_1,\ldots,e_m)$ belongs to $Z$ if the map
\[
 \Qp^m\longrightarrow N_t,\qquad
 (a_1,\ldots,a_m)\longmapsto\sum_i a_i e_i
\]
is injective at every geometric point.  The action of $\GL_m(\Zp)$ on
the tuple is the standard action; its action on functions is inverse
substitution.  This action must be distinguished from the dual-standard
action on the translation parameters of the lifting torsor.

\paragraph{The analytic joint frame space.}
Let $\mathcal X$ be the analytic realization over $E$ of the
joint frame tower $\{B_U\}$ in Proposition~\ref{prop:frame-torsor},
together with the lift tower \eqref{rings:root-torsor}, on the
exact-height locus $x\ne0$ of the open $A$-disc.  We take the
completed perfection of this finite-level system.  Forgetting the
connected marking gives a space $Y$, the perfected \'etale-basis
and lift tower on that locus, and $\mathcal X\to Y$ is a
$P_t$-torsor.  At each finite frame level, $\mathcal X/U$ is the
corresponding perfected analytic realization of $B_U$.
The algebraic union $\mathcal B=\colim_U B_U$ is still uncompleted,
with the common finite-stage and bounded-pole cochain convention.

The next lemma supplies coordinates for $Y$ by first retaining
the boundary.  We then use those coordinates to identify the
connected marking with a frame of a quotient bundle.

\begin{lemma}\label{support:tate-ball}
After scalar extension to $E$, the completed perfection of the full
etale-basis tower, including its boundary, is the product of $m$ open
Honda balls.  The exact-height locus is the open locus on which the
corresponding $m$ rational Tate sections are independent.  The
identification is compatible with $P_n$, changes of integral Tate basis,
and passage between finite frame levels.
\end{lemma}

\begin{proof}
Let $z_{i,l}$ be compatible lifted division coordinates, with
$[p]z_{i,l+1}=z_{i,l}$.  Proposition~\ref{prop:finite-rings} identifies
the finite lift algebras, including their infinitesimal fibers, with
the indicated root algebras.  Consequently, on perfectoid test rings
the lifts are unique after perfection.  The coordinates in the inverse
limit are
\[
 X_{i,l}=z_{i,l}^{p^{nl}},\qquad
 X_i=\lim_{l\to\infty}X_{i,l}.
\]
We check both convergence and the topology, since merely adjoining
roots to the generic ring would not determine the domain.

Put $I=(x_t,\ldots,x_{n-1})$ and
$\lambda=\max_{t\leq a<n}|x_a|<1$.  The full Honda congruence and the
factorization through $T^{p^t}$ in Section~\ref{sec:rings} give
\[
 [p](T)-T^{p^n}\in I T^{p^t}A[[T]],\qquad
 |[p]z-z^{p^n}|\leq\lambda |z|^{p^t}\quad (|z|<1).
\]
For $0<\rho<1$, define the closed level-$l$ domain and its unnormalized
radius by
\[
 B_l(\rho)=\left\{\max_i|z_{i,l}|^{p^{nl}}\leq\rho\right\},
 \qquad R_l=\rho^{p^{-nl}},\qquad l\geq1.
\]
For compatible division points, the characteristic-$p$ power identity
and the preceding error estimate imply
\[
 |X_{i,l+1}-X_{i,l}|
 \leq\lambda^{p^{nl}}|z_{i,l+1}|^{p^{nl+t}}
 \leq\lambda^{p^{nl}}.
\]
The proof of Proposition~\ref{prop:finite-rings} gives
$x_a\in(w_{1,1},\ldots,w_{m,1})$ in $C_1$, with
$w_{i,1}=z_{i,1}^{p^t}$.  Thus the level-one coordinate ideal, followed
by the increment estimate down from a point of $B_l(\rho)$, gives
\[
 \lambda\leq\max_i|z_{i,1}|^{p^t}
 \leq\max\{\rho^{p^{t-n}},\lambda^{p^t}\}.
\]
Indeed the increment estimate gives
$\max_i|X_{i,1}|\leq\max\{\rho,\lambda^{p^n}\}$, which is the second
inequality after taking the $p^{t-n}$ power.  Since $\lambda<1$ and
$p^t>1$, it follows that $\lambda\leq\rho^{p^{t-n}}$ uniformly on
$B_l(\rho)$.

These domains are preserved by the transition maps, which are also
surjective on them.  For $l\geq1$ one has
\[
 p^{-nl}(1-p^{t-n})<p^{t-n},\qquad
 \lambda\leq\rho^{p^{t-n}}
 <\rho^{p^{-nl}(1-p^{t-n})}=R_{l+1}^{p^n-p^t}.
\]
The error estimate therefore sends $B_{l+1}(\rho)$ into $B_l(\rho)$.
Conversely, finite lying-over for the frame transition supplies a
lifted valued point above any point of $B_l(\rho)$.  If one of its
coordinates $z$ had $|z|>R_{l+1}$, the strict inequality
$\lambda |z|^{p^t}<|z|^{p^n}$ would make the Honda term dominate, so
\[
 |[p]z|=|z|^{p^n}>R_l,
\]
contrary to the chosen point of $B_l(\rho)$.  Thus the lift belongs to
$B_{l+1}(\rho)$.

On these compatible domains the increment estimate is bounded by
$\rho^{p^{nl+t-n}}$, which tends to zero.  It proves uniform convergence
of the sequence defining $X_i$, with
\[
 |X_i-X_{i,l}|\leq\delta_l:=\rho^{p^{nl+t-n}},
 \qquad |X_i^{1/p^{nl}}-z_{i,l}|\leq\lambda.
\]
We verify the resulting map of completed coordinate rings on a
dense subalgebra.  Fix a nonzero polynomial $f$ over $E$ in
$X_1^{1/p^a},\ldots,X_m^{1/p^a}$, with one fixed denominator $p^a$.
For $nl\geq a$, its level-$l$ approximation is the polynomial
\[
 f(z_{1,l}^{p^{nl-a}},\ldots,z_{m,l}^{p^{nl-a}}).
\]
By Proposition~\ref{prop:finite-rings}, $B_l(\rho)$ is the closed
$z$-polydisc of radius $R_l$.  The displayed polynomial has exactly
the Gauss norm of $f$ at radius $\rho^{1/p^a}$; projection from the
inverse limit onto $B_l(\rho)$ is surjective, so its supremum norm
there is unchanged.  The coordinate errors are at most
$\delta_l^{1/p^a}$, which tends to zero.  Finite polynomial
substitution therefore converges uniformly to
$f(X_1^{1/p^a},\ldots,X_m^{1/p^a})$.  Once the error is smaller
than the nonzero Gauss norm of $f$, the ultrametric inequality makes
the two norms equal.  This proves that the map from the fractional
polynomial algebra is isometric and hence injective on its completion.

To prove density in the other direction,
\eqref{rings:transition-powers} expresses each coordinate from level
$l_0$ as an integral series in the $p^{t(l-l_0)}$-powers of the level-$l$
coordinates.  Substitution of $X_i^{1/p^{nl}}$ changes this series by at
most $\lambda^{p^{t(l-l_0)}}$, which tends uniformly to zero.  Thus the
$X_i$ and their roots topologically generate the inverse limit, proving
the asserted identification of perfectoid spaces.  Applying the estimates
to the addition law and to the finite-level change-of-frame formulas
proves equivariance.

A rational relation can be multiplied by a power of $p$ and then made
primitive over $\Zp$.  Its reduction at a sufficiently large division
level is a relation among the labelled basis sections.  The full-set
polynomial of Proposition~\ref{prop:finite-rings} says precisely that
such a relation occurs when the etale rank drops below $m$.  Conversely
a drop of that rank supplies a relation.  This proves the last
assertion, including at the boundary.
\end{proof}

We shall use the following curve facts with their relative meanings.
Standard positive bundles have the Honda universal covers as their
section sheaves
\cite[Corollary~5.1.2 and Proposition~5.1.6]{ScholzeWeinstein};
equivalently one may use
\cite[Theorems~15.2.3 and~15.2.5]{ScholzeWeinsteinBerkeley}.
Families with constant slope polygon become standard pro-etale locally,
and slope-zero bundles correspond to $\Qp$-local systems
\cite[Theorem~II.2.19 and Corollary~II.2.20]{FarguesScholze}.
The sheafified first cohomology of $\cO$ vanishes
\cite[Proposition~II.2.5(ii)]{FarguesScholze}.
The functor of section sheaves on positive bundles is fully faithful
over a varying perfectoid base: this follows from the fully faithful
$R\tau_*$ on perfect complexes and the positive-slope identification
with Banach--Colmez sheaves
\cite[Corollary~3.11]{AnschuetzLeBrasFourier}.
All maps to which we apply this assertion are $\Qp$-linear.

\begin{theorem}\label{thm:geometric-quotient}
Let $\mathcal X$ be the full perfected joint frame space over $E$.
There is a $P_n\times U_0$-equivariant description of $\mathcal X$ as
the space of exact sequences
\begin{equation}\label{support:framed-extension}
 0\longrightarrow\cO^m\longrightarrow V_n
 \longrightarrow V_t\longrightarrow0
\end{equation}
whose determinant comparison is a $\Zp^\times$-multiple of a fixed
comparison.  Forgetting the middle frame gives a $P_n$-torsor
$\mathcal X\to Z$.  Thus, for every compact open $U\subset U_0$,
\begin{equation}\label{support:quotient-stacks}
 [\mathcal X/U\,/P_n]\simeq[Z/U].
\end{equation}
On the quotient there is the universal sequence
\[
 0\longrightarrow L_m\otimes\cO\longrightarrow V_n'
 \longrightarrow V_t'\longrightarrow0,
\]
where $L_m$ has its integral etale lattice and the two positive bundles
have compact frame reductions.  Finally,
\begin{equation}\label{support:additive-presentation}
 N_t\simeq\mathbb A_C^{1,\diamond}/F,
\end{equation}
where $F/\Qp$ is unramified of degree $t$, acting by translations.
\end{theorem}

\begin{proof}
The universal-cover comparison sends the independent Tate sections
of Lemma~\ref{support:tate-ball} to a map $\cO^m\to V_n$.
Let its saturated image on a geometric fiber have rank $r\leq m$.
It is generically generated by sections, so has no negative
quotient and has degree at least zero.  Stability of $V_n$ bounds
its degree by $r/n<1$.  The degree is therefore zero, and all
its slopes are zero.  It is a trivial bundle of rank $r$, whose
rational section space has dimension $r$.  Independence forces
$r=m$.  The spanning determinant section has degree zero and no
zeros.  Thus the map is a subbundle inclusion on every fiber,
and hence relatively.  Every slope of its quotient is positive: a
nonpositive quotient would give a nonpositive quotient of $V_n$.
The quotient has rank $t$ and degree one.  By the classification of
bundles it is of type $V_t$, since two distinct positive summands
would already have total degree at least two.

Conversely, the middle term of an extension of $V_t$ by $\cO^m$ has
no negative slope: a negative quotient receives no map from either
end term.  A nonbasic middle term must consequently contain a
slope-zero quotient, since its total degree is one.  Such a quotient
is locally a trivial rational line, and it is equivalent to a rational
linear relation among the extension classes.  Thus the basic locus
is exactly $Z$.  These assertions are relative assertions after
pro-etale standardization and descend because the saturated slope
filtration is unique.

It remains to check that the connected marking gives the indicated
quotient frame, including its integral reduction.  Write
$\beta=\alpha^{-1}:\mathcal G_B^0\to\Gamma_t$ and put $q=p^t$.
On an affinoid perfectoid test over a bounded characteristic-$p$ chart,
write $[p]_{\mathcal G}(T)=a_qT^q+\cdots$ and
$\beta(T)=\sum_{j\geq1}b_jT^j$.  The coefficients of $[p]_{\mathcal G}$
are power bounded, with $a_q$ equal to $x$ times the fixed integral unit.
Comparing degree $qj$ in
$\beta([p]_{\mathcal G}(T))=\beta(T)^q$ gives a monic equation
\[
 b_j^q-a_q^j b_j=P_j(b_1,\ldots,b_{j-1}),
\]
where $P_j$ has power-bounded coefficients.  Induction gives
$|b_j|\leq1$ for every $j$; in particular $b_1^{q-1}=a_q$.
Thus on each closed subball $|z|\leq\rho<1$ the series converges
uniformly, with tail after degree $N$ bounded by $\rho^{N+1}$.
It defines a functorial continuous map on all topologically nilpotent
coordinates.  Applying it componentwise to compatible inverse-$[p]$
sequences gives a continuous map of universal-cover $v$-sheaves;
each coordinate uses a strict radius bound, and compact parameter
families use finite such covers.  The formal scalar identities give
$\Zp$-linearity, and the shift on the inverse limit gives $\Qp$-linearity.

For a Tate coordinate $z_l$ of order $p^l$, its image satisfies
$\beta(z_l)^{p^{tl}}=0$.  Characteristic-$p$ perfectoid tests are reduced,
so this image is zero.  The map therefore kills the rational span of
the Tate sections.  This torsion vanishing is only asserted on these
reduced analytic tests, not on arbitrary nilpotent Artin algebras.
Let $\mathcal Q=V_n/\cO^m$ be the positive quotient bundle above, and
write $\tau_*$ for the relative section-sheaf functor.  The relative
universal-cover comparison and the sheafified vanishing
$R^1\tau_*\cO=0$ give an exact sequence of topological
$\Qp$-module $v$-sheaves
\[
 0\longrightarrow\underline{\Qp}^{\,m}\longrightarrow\tau_*V_n
 \longrightarrow\tau_*\mathcal Q\longrightarrow0.
\]
Hence the map just constructed factors uniquely through
$\tau_*\mathcal Q\to\tau_*V_t$.  These positive section objects are
concentrated in degree zero, so relative full faithfulness produces an
actual bundle map $\overline\beta:\mathcal Q\to V_t$.  On each geometric fiber,
$b_1\ne0$ and a sufficiently small nonzero formal point has nonzero
image; division-point projection is surjective.  The bundle map is
therefore nonzero on every fiber.  A nonzero map between the stable
bundles of this type is an isomorphism, so the relative map is an
isomorphism.

We now show that this construction recovers every quotient frame
in the required determinant component.  Over the independent
Tate-section space $Y$, the connected markings form the $P_t$-torsor
$\mathcal X\to Y$.  The quotient bundles form the family
$\mathcal Q=V_n/\cO^m$ constructed above.  Their frames form the
$D_t^\times$-torsor
\[
 \mathscr F=\underline{\operatorname{Isom}}_Y(\mathcal Q,V_t).
\]
The map $\alpha\mapsto\overline\beta$ just constructed is
$P_t$-equivariant from $\mathcal X$ to $\mathscr F$.

The ordered trivial subbundle identifies $\det\mathcal Q$ with
$\det V_n$.  Fix a comparison $\kappa:\det V_n\simeq\det V_t$.
For $f\in\mathscr F$, put
\[
 \delta(f)=v_p\bigl(\det(f)/\kappa\bigr)\in\Z.
\]
Each fiber of this locally constant determinant map is a
$P_t$-torsor, because the kernel of
$v_p\operatorname{Nrd}:D_t^\times\to\Z$ is $P_t$.
In particular $\delta(\overline\beta)$ is unchanged by changing
the connected integral marking.  It descends to a locally constant
integer on $Y$; an integral change of Tate basis also leaves it
unchanged.

The space $Y$ is geometrically connected.  In the Honda-ball
coordinates of Lemma~\ref{support:tate-ball}, for each nonzero
rational row $v$ choose a nonzero integral multiple and let $f_v$
be the coordinate function of the corresponding Honda-linear
combination.  Its zero locus detects that rational relation.
At a point of $Y$, every $f_v$ is nonzero.  Pass to the completed
residue field of the point and take centered Gauss seminorms of
increasing radius less than one.  The Gauss expansion of each $f_v$
retains its nonzero constant term, so its norm stays nonzero along
the entire radius path.  This works separately for every $v$ and
requires no common lower bound for their norms.  Once the radius
contains the point's coordinates, the centered Gauss polydisc is
the origin Gauss polydisc.  Increasing its radius joins these
origin polydiscs.  Thus these paths remain in $Y$ and prove
connectedness on Berkovich points; rank-one maximalizations give
the same conclusion for locally constant functions on the adic
space.

Consequently $\delta(\overline\beta)$ has one value.  Rescale
$\kappa$ to make that value zero.  The map
$\mathcal X\to\mathscr F_0$ is now an equivariant map between
$P_t$-torsors over $Y$, hence an isomorphism.  This also gives the
inverse construction: a determinant-matched framed extension gives
its independent Tate sections in $Y$ and its quotient frame in
$\mathscr F_0$, which uniquely recover the connected marking.

Now the possible middle frames form a torsor under $D_n^\times$.
Their determinant valuations are the reduced-norm valuations, so
the determinant-matched frames form exactly the subgroup
$\cO_{D_n}^\times=P_n$.  This proves the torsor assertion and the
quotient equivalence.  This normalization also agrees with the
determinant condition in the two-tower description
\cite[Theorem~2.0.1]{BMR}: our finite connected and etale frames,
followed by the root lifting tower, are exactly the generic
trivialization there.  Indeed the finite lift algebras are the root
algebras of Proposition~\ref{prop:finite-rings}, and on perfectoid
tests their connected lifting torsors have unique sections.

For the last assertion use the curve with coefficient field $F$.
The untilt divisor gives
\[
 0\longrightarrow\cO_F(-1)\longrightarrow\cO_F
 \longrightarrow i_*\cO_C\longrightarrow0.
\]
The section sheaf of $\cO_F$ is $F$, its first cohomology sheaf
vanishes, and $\cO_F(-1)$ has no sections.  Finite pushforward to
the $\Qp$-curve identifies $\cO_F(-1)$ with $V_t^\vee$ by rank,
degree, and the cyclic unramified Frobenius description.  The
cohomology sequence is therefore $0\to F\to\mathbb A_C^{1,\diamond}
\to N_t\to0$, proving \eqref{support:additive-presentation} as a
sequence of sheaves on arbitrary perfectoid tests.
\end{proof}

\subsection{The integral support convention}

All supports below are computed on bounded closed subdomains lying
overconvergently inside an open chart.  Equivalently, one takes the
filtered colimit of the fibers of restriction to the complements of
such subdomains.  Inner and outer radii always have a strict gap.
This convention agrees with compactly supported pushforward for the
partially proper spaces used here.  On quotient presentations we
first choose such supports upstairs and then descend their complexes.

We record the localization facts needed to use this convention with
$\mathscr A$.  The integral structure sheaf modulo $p$ is an etale
coefficient.  For a quasi-pro-etale map $f:X\to Y$, its pullback is
$f^*(\mathcal O_Y^+/p)\simeq\mathcal O_X^+/p$
\cite[Proposition~2.1.2 and the proof of Lemma~2.1.3]{PignonYwanne}.
For a cofiltered spatial system with coefficient pulled back from
one stage, its cohomology commutes with the inverse limit of spaces
\cite[Proposition~14.9]{ScholzeDiamonds}.  Quasicompact injections
are quasi-pro-etale \cite[Corollary~10.6]{ScholzeDiamonds}.
Consequently shrinking rational tubes to a generalizing closed
locus computes restriction to that locus.  The open--closed
localization triangle for this coefficient, also in its solid almost
form, is
\begin{equation}\label{support:localization}
 j_!j^*K\longrightarrow K\longrightarrow i_*i^*K
\end{equation}
for the specializing open complements occurring here
\cite[Corollary~3.2.4 and Proposition~4.3.1]{PignonYwanne}.
In particular this use of localization does not assert coherent
purity across an arbitrary characteristic-$p$ boundary.

The compact group calculations will use the following finite
resolution model.  It applies both to the support complexes here
and to the geometric constant-cochain comparison later.

\begin{lemma}[Finite projective calculation of solid rows]
\label{lem:solid-rows}
Let $K$ be a compact $p$-adic analytic group without $p$-torsion, and put
$R_K=\Fp[[K]]$.  Let $\mathcal V$ be a solid $\Fp$-module with continuous
condensed $K$-action.  Choose a bounded resolution $Q_\bullet\to\Fp$ by
finitely generated projective profinite $R_K$-modules.  Then continuous
solid $K$-cohomology is computed by
\begin{equation}
 \underline{\Hom}_{\underline{R_K}}
           (\underline{Q_\bullet},\mathcal V).
 \label{eq:const-solid-resolution}
\end{equation}
Evaluation of this complex at the one-point profinite set is exact and
each of its terms is a retract of a finite direct sum of $\mathcal V(*)$.
In particular, if a commuting endomorphism $u$ acts as a scalar $c$ on
$\mathcal V(*)$, it acts as $c$ on the point-valued cohomology of
\eqref{eq:const-solid-resolution}.

For a derived solid coefficient $\mathcal A$, this gives a natural
spectral sequence
\begin{equation}
 H^r_{\mathrm{solid}}(K;H^s(\mathcal A))(*)
      \ \Longrightarrow\
 H^{r+s}(\RGamma_{\mathrm{solid}}(K;\mathcal A))(*).
 \label{eq:const-solid-ss}
\end{equation}
The preceding scalar assertion applies to every row.
\end{lemma}

\begin{proof}
The existence of the chosen bounded resolution follows from
Noetherianity and finite global dimension of $\Fp[[K]]$ when $K$ has
no element of order $p$ \cite[Proposition~3.3]{ArdakovWadsley}; finite
generation of its terms follows by successively resolving kernels.
The condensation functor from profinite $R_K$-modules is exact and preserves
projectives by \cite[Theorems~3.2 and~3.14(i)]{Tang}.  Moreover, solid
$\Fp$-modules with commuting condensed $K$-action are precisely solid
$\underline{\Fp[[K]]}$-modules by
\cite[Proposition~3.21]{Tang}.  In this assertion the solidification of
the condensed group ring is the completed group ring, including its
multiplication.  Thus condensation of $Q_\bullet$ is a projective
resolution in the required coefficient category, proving
\eqref{eq:const-solid-resolution}.

Write $Q_i$ as a summand of $R_K^{N_i}$.  Internal Hom from this finite
free module is $\mathcal V^{N_i}$, and the idempotent defining $Q_i$
cuts out the corresponding summand.  Evaluation at $*$ preserves finite
sums and these idempotents, and is exact: the point is an extremally
disconnected projective object of the condensed site.  A scalar on
$\mathcal V(*)$ therefore stays that scalar on every term and on its
cohomology.  No assertion that point evaluation is conservative on solid
modules is involved.  Filtering $\mathcal A$ by its cohomology objects
gives \eqref{eq:const-solid-ss}; the finite length of $Q_\bullet$
controls its group-cohomological direction.
\end{proof}

Here is the comparison with geometric descent.  For a locally spatial
diamond $Y'$ and an extremally disconnected profinite set $T'$, form
the finite-coefficient geometric cochains on $Y'\times T'$.
An affinoid-perfectoid hypercover and the Artin--Schreier complex compute
these cochains.  On a characteristic-$p$ perfectoid affinoid the analytic
ring with its topology is a complete linearly topologized $\Fp$-module,
hence an inverse limit of discrete $\Fp$-modules.  On its product with
$T'$ the ring is the corresponding continuous-function module.  These
are solid.  Kernels, extensions, and derived limits retain solidity, so
the Artin--Schreier complexes and their hypercover totalizations are
derived solid.  The derived closure assertion for profinite coefficient
rings is also recorded in \cite[Theorem~4.4(ii)]{Tang}.  Evaluation at
$*$ gives the usual finite-coefficient geometric \'{e}tale cochains.
The nerve of a compact group action is the same parameter construction.

To compare that nerve with \eqref{eq:const-solid-resolution}, solidify
the free bar resolution.  For a profinite set $T'$, the derived
solidification of the free condensed $\Fp$-module on $T'$ is
$\Fp[[T']]$ in degree zero: resolve $T'$ by extremally disconnected
profinite sets and dualize the resulting augmented complex to locally
constant $\Fp$-functions, which are acyclic on a profinite space.
The solidified bar terms are therefore the usual free profinite
completed-group-ring modules.  The augmented bar is still exact,
with its usual continuous contraction after forgetting the group
action.  It is a projective resolution by Tang's
projective-preservation theorem, and comparison with $Q_\bullet$
gives the asserted equality of descent constructions.  This argument
uses analytic rings to construct the coefficient category; it does
not replace a topological geometric coefficient by a discrete tensor
product.

\begin{lemma}\label{support:parameters}
These localization and support computations are compatible with
compact profinite parameters and compact frame torsors.  They apply
to the rational-rank loci in $N_t^m$ and to their flag incidences.
On an exact-rank locus the smallest rational plane depends
continuously on the point.
\end{lemma}

\begin{proof}
\emph{Integral continuity and compact parameters.}
On an affinoid perfectoid, integral sections modulo $\varpi$ compute
$\mathscr A$ up to almost isomorphism.  This is affinoid almost
acyclicity and its v-descent form
\cite[Proposition~8.8]{ScholzeDiamonds}.  Completing a union of
affinoid perfectoid rings does not change its reduced integral
sections almost: for every $\epsilon>0$, quasicompactness makes a
bound on the limit an eventual bound with norm loss at most
$|\varpi|^{-\epsilon}$.  The same estimate applies to the
$\varpi$-multiples.  Finite Cech diagrams preserve these estimates.
Alternatively one may pull a tube system to a strictly totally
disconnected cover; its terms and nerve are then computed by
rational subsets and their buffered limits.  This verifies that the
coefficient is pulled back in the continuity assertion just stated.

For a compact profinite parameter, reduction modulo a fixed
positive-valuation cutoff makes continuous sections locally
constant.  Uniform approximation on finitely many clopen pieces
therefore reduces the calculation to the one without parameters.
The same argument applies on the nerve of a compact profinite
torsor.  Such a torsor is proper, and inverse images of compact
bounds and their compact-group translates admit common compact
bounds.

\emph{Compact groups and the support colimit.}
We justify passage through the support colimit in the stable almost
localization of derived solid $A_0$-modules, where
$A_0=\cO_E/\varpi$ and $A_0^a=R$.  For a buffered support $K$, let
$\mathscr C_K$ be its supported parameter object.  The cutoff
continuous-function modules above are solid $\Fp$-modules; their derived
Cech limits and supported fibers give $\mathscr C_K$ in this category.
The compact-support object is $\colim_K\mathscr C_K$ there.
Every compact frame group $G$ used here is $p$-adic analytic without
$p$-torsion, since $p>n+1$.  Choose a bounded resolution of $\Fp$ by
finitely generated projective profinite $\Fp[[G]]$-modules
\cite[Proposition~3.3]{ArdakovWadsley}.  The continuous solid bar
calculation agrees with this resolution
\cite[Theorems~3.2 and~3.14(i), Proposition~3.21]{Tang}; the completed-bar
comparison is detailed following Lemma~\ref{lem:solid-rows}.
Derived solid scalar extension to $R$ preserves its augmentation and
makes its terms retracts of finite free modules over the $R$-linear
action algebra.  Thus invariants are computed by a bounded complex
of retracts of finite sums of the coefficient object.  Forming both
computations inside the almost category gives
\[
 \colim_K\RGamma(G,\mathscr C_K)
   \xrightarrow{\ \sim\ }
 \RGamma\bigl(G,\colim_K\mathscr C_K\bigr).
\]
Only the group resolution is bounded; no amplitude bound on
$\mathscr C_K$ is required.

For the two torsors in \eqref{support:quotient-stacks}, choose
common cofinal buffered supports upstairs invariant under both groups,
using properness.  At a fixed support, the supported fiber commutes with Cech
descent, and the two orders of descent give the common double nerve
$\mathcal X\times P_n^a\times U^b$.  Replacing its two group directions
by their bounded resolutions gives a finite double complex whose terms
are retracts of finite sums.  The support colimit therefore commutes with
both directions.  This proves that the torsor nerves compute compact
supports and that the two presentations agree.  The comparison remains
the canonical double-nerve map; its support, cup, and trace maps are
unchanged by this verification.

\emph{Rational-rank incidence.}
On bounded
affine lifts of $N_t^m$, cover the compact rational Grassmannian
by finitely many clopen charts admitting matrix frames.  In such a
chart, incidence with a plane of dimension $j$ has equations
\[
 A(L)z=b,\qquad b\in F^{m-j}.
\]
The entries of $A(L)$ have a common bound, say $C_0$, and $z$ has
some bound $M_0$.  Therefore every possible $b$ lies in the compact
ball $|b|\leq C_0M_0$ in $F^{m-j}$.  Thus all relevant relations
are included in a closed analytic incidence over compact
parameters.  Projection along these parameters is proper.  Its
image is closed and generalizing: the same equation and the same
witness survive generization.  Equivalently one may use the
generalizing property of maps of locally spatial diamonds
\cite[Proposition~11.19(iv), Definition~18.1, and Remark~18.2]{ScholzeDiamonds}.
Its open complement is partially proper by the valuative
criterion \cite[Definition~18.4]{ScholzeDiamonds}.

Choose a finite mesh in the compact parameter chart, bound the
values of its relation functions, and let those bounds tend to
zero.  The resulting rational tubes have this incidence image as
their closed limit.  Strict radius gaps exclude additional
valuations at an outer boundary.  The construction is compatible
with every quasicompact test and with refinements of the mesh.
On the stratum of actual rank $j$, the smallest plane is unique.
The proper incidence projection with this unique fiber has a
continuous inverse there.  This is precisely the continuity used
when contracting the flag choices on an exact-rank stratum.
\end{proof}

\begin{lemma}\label{support:affine}
For the diamond of the untilted affine space one has
\[
 \RGamma_c(\mathbb A_C^{j,\diamond},\mathscr A)
 \simeq R(-j)[-2j].
\]
Translations and invertible linear changes act trivially on this
line, apart from its displayed arithmetic Tate twist.  The
identification is compatible with compact profinite parameters.
\end{lemma}

\begin{proof}
Take $\mathbb P_C^j$ with hyperplane complement
$\mathbb A_C^j$.  Proper primitive comparison
\cite[Theorem~5.1]{ScholzeHodge} and
\eqref{support:localization} identify the compact-support complex
with the fiber of the projective-space restriction map.  Restriction
preserves the powers of the hyperplane class, so cancels the even
degrees from $0$ through $2j-2$ and leaves $R(-j)$ in degree $2j$.
The same conclusion follows directly from compact-support primitive
comparison for Zariski opens of smooth proper spaces
\cite[Corollary~5.3.2 and its proof]{PignonYwanne}, with
algebraic--analytic proper-support compatibility
\cite[Proposition~3.16]{BhattHansen}.
The top affine class is invariant under affine linear changes,
and the comparison and localization maps are natural.
Lemma~\ref{support:parameters} proves the parameter assertion.

For clarity, the boundary germ involved here is indeed the
coefficient at a point.  Shrinking quasicompact discs to that point
and applying spatial continuity gives $R$ in degree zero.  One can
also see the vanishing directly on a Kummer annulus: a nonconstant
fractional monomial of exponent $u=bp^v$, $p\nmid b$, has
difference operator of norm $|\epsilon-1|^{p^v}$.  Shrinking both
Laurent bounds by a ratio at most $|\epsilon-1|$ makes division by
this operator integral, since $p^v\leq|u|_{\mathbb R}$.  This
contracts the nonconstant complex.  The remaining degree-one
Kummer class vanishes on a smaller disc by the binomial-series
root.  This argument also explains why the same assertion is not
being made about compact supports of a perfected open unit disc.
\end{proof}

\begin{lemma}\label{support:translation-quotient}
For $j\geq0$,
\begin{equation}\label{support:ambient-cohomology}
 \RGamma_c(N_t^j,\mathscr A)
 \simeq R(-j)\otimes\operatorname{or}_{F^j}^{-1}[-(t+2)j].
\end{equation}
Here $\operatorname{or}_{F^j}^{-1}$ is the top continuous
cohomology orientation of a lattice in the $\Qp$-space $F^j$.
On $\GL_j(\Zp)$ its character is $(\det)^{-t}$, reduced to the
coefficient ring.  The identification retains the commuting frame
actions and compact parameters.
\end{lemma}

\begin{proof}
Let $\Lambda_s=p^{-s}\cO_F^j$.  Choose radii
$|p|^{-r}<\rho_r<|p|^{-(r+1)}$ and let $D_r$ be the open
polydisc of radius $\rho_r$, with buffered compact supports inside.
Then $F^j\cap D_r=\Lambda_r$, and translates by elements outside
$\Lambda_r$ are disjoint from $D_r$.  For a fixed
lattice, compact-group descent commutes with exhaustion by
invariant supports.  Lemma~\ref{support:affine} consequently gives
\[
 \colim_r\RGamma_c(D_r/\Lambda_s,\mathscr A)
 \simeq
 \RGamma(\Lambda_s,R(-j))[-2j].
\]
The quotients $D_r/\Lambda_r$, with radii chosen to contain exactly
the prescribed translation lattice, embed openly into $N_t^j$
and exhaust it.  More explicitly, use pairs $(s,r)$ with $r\geq s$.
Increasing $r$ gives extension on the disc, while increasing $s$
gives corestriction for the finite cover between lattice quotients.
These are the maps induced by the open embedding into the next
quotient chart: on the cover, they sum the disjoint translates.
The diagonal is cofinal, because a bounded compact support and
any finite compact family of translates fit in a later disc.
Thus one may first exhaust the affine space at fixed $s$ and then
take the lattice colimit.

Put $h=tj$.  The continuous Koszul complex for the abelian lattice
gives
\[
 H^a(\Lambda_s,R)=
 \bigwedge^a\Hom_{\mathrm{cts}}(\Lambda_s,\Fp)\otimes_{\Fp}R.
\]
In homothety frames the corestriction from $\Lambda_s$ to
$\Lambda_{s+1}$ multiplies degree $a$ by $p^{h-a}$.  This follows
either from the exterior Koszul comparison or by duality with
restriction in complementary degree.  Hence the colimit kills
every degree except $h$, where it preserves the top orientation.
Together with the affine degree $2j$ this proves
\eqref{support:ambient-cohomology}.  Under inverse substitution an
automorphism acts on the top lattice class by the inverse of its
$\Qp$-determinant.  A matrix on $j$ rows acts on $F^j$ with
determinant $(\det)^t$, proving the stated character.  The finite
Koszul and compact-support maps are functorial, so the argument
retains all commuting actions.

For the commuting $P_t$-action this functoriality can be expressed
without choosing a $P_t$-invariant affine chart.  The constant-$\Fp$
primitive comparison on the affine pieces, followed by the same
lattice corestriction colimit, gives the canonical equivalence
\[
 \RGamma_c(N_t^j,\Fp)\otimes_{\Fp}R
       \xrightarrow{\ \sim\ }\RGamma_c(N_t^j,\mathscr A).
\]
The source has a single one-dimensional $\Fp$-cohomology group,
in degree $(t+2)j$.  The comparison is natural for automorphisms
of $N_t^j$ and for compact parameter families.  Its $P_t$-action
is consequently a smooth finite character, as required when
taking compact-group invariants below.
\end{proof}


\subsection{The rational flag resolution}

To recover the compact supports of $Z$ from those of $N_t^m$, we resolve
the complement of $Z$, consisting of dependent tuples, by rational supporting planes allowed
to vary over compact Grassmannians. Flags organize the incidence among these
plane families. The augmented incidence complex below represents extension
by zero from $Z$, so the next definitions record the orientation line and
flag cohomology associated to each smallest plane.

For $m=1$ the complement is only the zero extension.  The support
triangle is
\[
 \RGamma_c(N_t\setminus\{0\},\mathscr A)
 \longrightarrow\RGamma_c(N_t,\mathscr A)
 \longrightarrow R.
\]
Lemma~\ref{support:translation-quotient} places the ambient line
$R(-1)\otimes\operatorname{or}_F^{-1}$ in degree $t+2$.
The triangle therefore gives a boundary copy of $R$ in degree one
and that ambient line in degree $t+2$.  For general $m$, the
proper rational planes replace the single zero plane, and their
flags account for the additional incidence degrees.

For a rational plane $W\subset\Qp^m$ of dimension $j$, write
$N_t(W)=W\otimes_{\Qp}N_t$ and
\[
 \mathfrak o_t(W)=
 \bigwedge^{tj}_{\Fp}
 \Hom_{\mathrm{cts}}((W\cap\Zp^m)\otimes_{\Zp}\cO_F,\Fp)
 \otimes_{\Fp}R.
\]
This is the line in Lemma~\ref{support:translation-quotient}; its
plane-row character is $(\det)^{-t}$.  Its natural commuting $P_t$
action is retained.  We put $\mathfrak o_t(0)=R$.

If $Q$ has rational dimension $s$, let
$\Delta_s=\{1,\ldots,s-1\}$, and denote by
$\operatorname{Fl}_I(Q)$ the compact space of rational flags with
dimension set $I\subset\Delta_s$.  The empty flag space is a point.
For a coefficient module $M$, consider the augmented complex
\begin{equation}\label{support:flag-complex-definition}
 \mathcal B_Q^a(M)=
 \bigoplus_{\substack{I\subset\Delta_s\\|I|=a-1}}
 \operatorname{LC}(\operatorname{Fl}_I(Q),M),
 \qquad 1\leq a\leq s.
\end{equation}
Its differential adding $i\notin I$ is pullback along the flag
projection, with sign $(-1)^{|\{u\in I:u<i\}|}$.
Define
\[
 \operatorname{St}(Q;M)=
 \frac{\operatorname{LC}(\operatorname{Fl}_{\Delta_s}(Q),M)}
 {\displaystyle\sum_{i\in\Delta_s}
   \operatorname{LC}(\operatorname{Fl}_{\Delta_s\setminus\{i\}}(Q),M)},
\]
where the maps in the denominator are pullbacks.
For $s=1$ this means $\operatorname{St}(Q;M)=M$.

\begin{lemma}\label{support:flag-complex}
The complex \eqref{support:flag-complex-definition} has cohomology
only in degree $s$, where it is $\operatorname{St}(Q;M)$.
The assertion is natural under linear isomorphisms and for locally
constant families over compact rational Grassmannians.
\end{lemma}

\begin{proof}
Choose a basis and an upper triangular Borel.  For each partial
flag type, filter by the decreasing Bruhat opens consisting of
cells of length at least $b$.  Restriction gives the associated
graded as compactly supported locally constant functions on the
cells of length $b$.  These restriction maps are surjective:
compactly supported locally constant functions on a closed stratum
extend after a finite refinement by compact open neighborhoods.

A flag projection sends a refined cell indexed by $v$ to that
indexed by the shortest representative $w$ of its parabolic
coset.  Its length can only decrease.  If the lengths agree, then
$v=w$ and the projection is an isomorphism on the common
unipotent cell coordinates.  Hence, on the associated graded,
only these isomorphisms contribute.

For a permutation $w$, the resulting summand is the signed
Boolean complex on the subsets $I$ containing its right descent
set $D_R(w)$, with the same cell coefficient in degree $|I|+1$.
If $D_R(w)\ne\Delta_s$, toggling any element of
$\Delta_s\setminus D_R(w)$, with the incidence sign, contracts
this complex.  Only the longest permutation survives, in degree
$s$.  The filtration is finite, so this proves vanishing in all
other degrees.  Its top cokernel is the displayed Steinberg
quotient.

Over a Grassmannian choose a finite disjoint clopen refinement
with quotient frames.  Locally constant sections commute with
this calculation: they are filtered colimits over finite clopen
partitions of finite products.  The constructions are natural,
so the local identifications glue independently of the frames.
\end{proof}

For $0\leq j<m$, let $\mathcal V_j$ be the locally constant
sections over $\operatorname{Gr}_j(\Qp^m)$ of the family
\[
 W\longmapsto
 \mathfrak o_t(W)(-j)\otimes_R\operatorname{St}(\Qp^m/W;R).
\]
This is locally constant induction from the compact parabolic
stabilizing a $j$-plane.  Its first block has orientation character
$-t$ and its remaining block has character zero before taking the
flag complex.  Put
$\mathcal V_m=\mathfrak o_t(\Qp^m)(-m)$.

\begin{proposition}\label{prop:compact-support}
The compactly supported cohomology of $Z$ is
\[
 H_c^a(Z,\mathscr A)=
 \begin{cases}
  \mathcal V_j,&a=(t+2)j+(m-j),\quad 0\leq j<m,\\
  \mathcal V_m,&a=(t+2)m,\\
  0,&\text{otherwise}.
 \end{cases}
\]
The map to $\RGamma_c(N_t^m,\mathscr A)$ identifies the last
cohomology group with the ambient line.  Scalar matrices in
$1+p\Zp$ act trivially on these cohomology groups.  The statement
is equivariant for the commuting compact frame groups and retains
the actual finite incidence complex, without choosing an
equivariant splitting of that complex.
\end{proposition}

\begin{proof}
For $a\geq1$, let $\mathcal I_a$ parametrize
\[
 (z,W_0<\cdots<W_{a-1}<\Qp^m),\qquad z\in N_t(W_0).
\]
The maps $\pi_a:\mathcal I_a\to N_t^m$ are proper by
Lemma~\ref{support:parameters}.  With the ambient term in degree
zero, form the alternating incidence complex
\begin{equation}\label{support:incidence-resolution}
 \mathscr A_{N_t^m}\longrightarrow
 \pi_{1*}\mathscr A_{\mathcal I_1}\longrightarrow\cdots
 \longrightarrow\pi_{m*}\mathscr A_{\mathcal I_m}.
\end{equation}
Forgetting the smallest plane uses restriction from a larger
plane fiber to a smaller one; the other face maps forget a flag
step.

This complex represents extension by zero from $Z$.  An
independent point has no proper containing plane.  At a dependent
point, the containing-plane poset has an initial object, its
smallest plane, so adjoining this plane contracts the augmented
nerve.  This is a contraction of the supported calculation, not
only a calculation of geometric-point ranks: filter by actual
rank and use the continuous smallest-plane map of
Lemma~\ref{support:parameters}.  On its finite clopen matrix
charts the contraction acts on the same tube and parameter
complexes.  Coefficient localization then glues the contraction.

Apply compactly supported cohomology to
\eqref{support:incidence-resolution}.  First compute the
plane-fiber direction by Lemma~\ref{support:translation-quotient}.
For a fixed smallest plane of dimension $j$, the remaining
complex is \eqref{support:flag-complex-definition} for the
quotient of dimension $m-j$, tensored with
$\mathfrak o_t(W)(-j)$.  Its term with no further plane lies in
incidence degree one.  The full incidence sign is the negative
of the flag sign; multiplication in flag degree $r+1$ by
$(-1)^r$ identifies the two conventions.  Lemma~\ref{support:flag-complex}
therefore gives the entries
\[
 E_2^{m-j,(t+2)j}=\mathcal V_j\quad(0\leq j<m),
 \qquad E_2^{0,(t+2)m}=\mathcal V_m.
\]
There are no further differentials.  A differential to an entry
with smaller plane dimension $j'<j$ would require
\[
 r=j-j',\qquad r-1=(t+2)(j-j'),
\]
which is impossible.  The total degrees are distinct, proving
the formula.  Projection onto the ambient term is the map induced
by extension of supports along $Z\subset N_t^m$, so its top edge
has the asserted identification.

A scalar fixes all rational flags.  On a plane of dimension $j$
it acts by $\overline a^{-tj}$ on the orientation, and trivially
on the affine Tate line.  This is one for $a\in1+p\Zp$.
The assertion follows on every displayed cohomology group.
All maps used in the construction are pullbacks, restrictions,
extensions of supports, or the lattice corestrictions described
above.  Thus the comparison retains the commuting actions.
\end{proof}

\begin{corollary}\label{support:finite-invariants}
For a compact open $U\subset U_0$ and a fixed finite character
twist, the groups
\[
 H^a\bigl(U,\RGamma_c(Z,\mathscr A)\otimes\chi\bigr)
\]
vanish outside a finite range and have bounded almost length over
$R$.  In particular, if such a group is $V\otimes_kR$ for a
$k$-vector space $V$, then $V$ is finite dimensional.
\end{corollary}

\begin{proof}
Use the finite flag resolutions rather than replace them by an
infinite representation without its resolution.  Shapiro reduces
the $\GL_m$-direction to compact parabolic groups with a line
coefficient.  Their intersection with an open subgroup has
finitely many orbits on each compact flag space.  The commuting
$P_t$-direction is compact as well.  All these groups have finite
$p$-cohomological dimension and a bounded resolution with finite
projective terms: they are compact $p$-adic analytic groups
without $p$-torsion, since $p>n+1$.
Their line coefficients and the finite twists are smooth; after
shrinking an open subgroup they are trivial.  Thus each computing
complex has finitely many finite almost free coefficient terms.
The same is true of the finite flag and group filtrations.

For the last assertion normalize the length of $R$ to one.
Subquotients and finite extensions of $R^b$ have length bounded
by the sum of the corresponding integers $b$; the assertion
also holds after almostification, by taking arbitrarily small
valuation losses.  A direct sum of $r$ copies of $R$ has length
$r$.  Every finite-dimensional subspace of $V$ therefore has
bounded dimension.  Faithfulness of the scalar extension implies
that $V$ is finite dimensional.
\end{proof}

\subsection{Principal parts and the positive Frobenius system}

We now return to a sufficiently deep finite joint frame level $U$.
Let $P=x^{-b}L$ be a finite free compact $C_1$-lattice in $B_U$,
with the following properties.  It is stable under the stationary
transfer $S$, every fixed larger pole lattice is carried into $P$
by a sufficiently large iterate, and the transpose on
\[
 M=\Hom_{C_1}(P,\omega_{C_1}),\qquad
 \omega_{C_1}=\Omega^m_{C_1/k,\mathrm{cts}},
\]
is positive Frobenius in the finite volume frame $\eta$.
More precisely, after the divisible choice of root step in
Lemma~\ref{lem:cartier-transfer}, write this step as $Q$-Frobenius.
Its matrix on $M$ has entries in $xC_1$.  In the etale trace and
volume identification on $x\ne0$, the basis functions of $M$
are integral over $C_1$.  These lattice choices exist by
Lemma~\ref{lem:compact:stable-lattices}; that lemma uses only the
finite-level rings and denominator bounds, not a cohomological
finiteness assertion.  The integrality condition can always be
arranged simultaneously with positivity: a sufficiently high
power of $x$ clears the coefficients of the finitely many monic
equations of the basis functions, while the Frobenius matrix gains
$(Q-1)b$ powers of $x$.

The continuous residue pairing identifies
\begin{equation}\label{support:residue-dual}
 P^\vee=\Hom_{\mathrm{cts}}(P,k)
 \simeq H^m_{\mathfrak m}(M),
 \qquad \mathfrak m=(w_1,\ldots,w_m).
\end{equation}
On a free basis it pairs power series with finite sums of monomials
negative in every variable, by the coefficient of
$w_1^{-1}\cdots w_m^{-1}dw_1\wedge\cdots\wedge dw_m$.
This is an equivariant pairing: the top differential includes the
Jacobian in the change-of-variables formula for the residue.

We specify the root coefficient system before passing to analytic
supports.  Choose a free basis $e_1,\ldots,e_N$ of $M$.  On $x\ne0$,
the trace and volume frame writes $e_i=f_i\eta$, with $f_i$ a
finite-cover function integral over $C_1$.  The actual Frobenius
transpose $F=S^\vee$ has the form
\[
 F(e_i)=f_i^Q\eta=\sum_j a_{ji}e_j,
 \qquad a_{ji}\in xC_1.
\]
Put $C^{(r)}=C_1^{1/Q^r}$ and let
$M_r=\bigoplus_i C^{(r)} e_i^{(r)}$ be free on the displayed symbols.
Only on the puncture do we interpret $e_i^{(r)}$ as
$f_i^{1/Q^r}\eta$; the volume line is retained once and is not rooted.
Define the transition, linear over $C^{(r)}\subset C^{(r+1)}$, by
\begin{equation}\label{support:root-coefficient-transition}
 e_i^{(r)}\longmapsto
       \sum_j a_{ji}^{1/Q^{r+1}}e_j^{(r+1)}.
\end{equation}
Taking $Q^{r+1}$-th roots of the displayed function identity verifies
that this is the actual inclusion on the puncture.  Its entries are
integral on the ambient rooted disc and are divisible by
$x^{1/Q^{r+1}}$.

Power renaming gives a $k$-linear isomorphism
$\theta_r:M_r\to M$, sending $c e_i^{(r)}$ to $c^{Q^r}e_i$.
It satisfies $\theta_{r+1}\iota_r=F\theta_r$, where $\iota_r$
is \eqref{support:root-coefficient-transition}.  The same identity
holds on the Cech fractions computing local cohomology.  Thus the
stationary transition on principal parts is exactly $F$.

On $\mathbb D^{m,\mathrm{perf}}$, let
$\mathscr M_r=\bigoplus_i\mathscr A e_i^{(r)}$, with the integral
analytic transition matrix \eqref{support:root-coefficient-transition}.
This defines the free ambient extension without assuming that the
cover functions themselves extend across $x=0$.  The maps $\theta_r$
are power renamings of the algebraic $k$-coefficients; after scalar
extension their comparisons fix $E$ and $R$, rather than applying
absolute Frobenius to those scalars.  The finite torsor actions and
morphisms of Section~\ref{sec:rings} commute with this system.

\begin{lemma}\label{support:principal-parts}
On the perfected open $m$-disc, compact supports of the free
integral lattice coefficient are computed by principal parts in
degree $m$.  For the root system of $M$, the resulting comparison is
\begin{equation}\label{support:principal-parts-comparison}
 \colim_{F} H^m_{\mathfrak m}(M)\otimes_kR[-m]
 \simeq \RGamma_c\left(\mathbb D^{m,\mathrm{perf}},
                    \colim_r\mathscr M_r\right),
\end{equation}
where $\mathscr M_r$ and its transitions are defined above, and
$F=S^\vee$ under \eqref{support:residue-dual}.
The map is natural for lattice
morphisms, integral coordinate changes, and continuous $P_n$
parameters.
\end{lemma}

\begin{proof}
First fix a coefficient stage.  Write $K_a$ for the closed
polydisc of radius $a$, $B_s$ for the outer closed polydisc of
radius $s<1$, and
$E_{b,s}=\bigcup_i\{b\leq|w_i|\leq s\}\subset B_s$ for its
Laurent exterior.  For $0<a<b<c<s<1$, put
\[
 \mathscr C_a=\RGamma_{K_a}(\mathbb D^{m,\mathrm{perf}},\mathscr M_r),
 \qquad
 \mathscr L_{b,s}=\fib\bigl(\RGamma(B_s,\mathscr M_r)
                    \to\RGamma(E_{b,s},\mathscr M_r)\bigr).
\]
The inclusions
$B_s\setminus K_c\subset E_{b,s}\subset B_s\setminus K_a$
give, by restriction and excision at the two ends, maps
\begin{equation}\label{support:buffer-sandwich}
 \mathscr C_a\longrightarrow\mathscr L_{b,s}\longrightarrow\mathscr C_c
\end{equation}
whose composite is the canonical enlargement of supports.

Here is the direction of the maps in the buffered system.
Index it by triples $(a,b,s)$ with $a<b<s$, and let a later
triple $(a',b',s')$ satisfy $s<a'$.  Choose $b<c<s$ and use
$\mathscr L_{b,s}\to\mathscr C_c\to\mathscr C_{a'}
\to\mathscr L_{b',s'}$ as its transition.
Comparing two choices with a larger intermediate $c$ shows
independence of that choice; \eqref{support:buffer-sandwich}
also shows that these transitions compose.  The genuine support
system and this buffered system interleave cofinally, so their
colimits agree.  Outer-domain comparisons enter through excision;
the forward maps enlarge supports.

The intersections in the finite cover of $E_{b,s}$ are
\[
 U_I=\{b\leq|w_i|\leq s\ (i\in I),\quad
            |w_j|\leq s\ (j\notin I)\},
 \qquad \varnothing\ne I\subseteq\{1,\ldots,m\}.
\]
Affinoid almost acyclicity computes $\mathscr L_{b,s}$ by the augmented
Cech complex of integral Laurent coefficients modulo $\varpi$.

For an exponent vector $\alpha$, a monomial $w^\alpha$ occurs on
$U_I$ exactly when its negative exponents have indices in $I$.
Wherever it occurs its Gauss weight is
\[
 b^{\sum_{\alpha_i<0}\alpha_i}
 s^{\sum_{\alpha_i\geq0}\alpha_i}.
\]
Thus the monomial Cech contraction preserves the norm: it removes
every exponent vector except those negative in all coordinates,
which remain in degree $m$ of the restriction fiber.  On a fixed
buffered chart the weighted coefficients of a convergent Laurent
series tend to zero.  Modulo any fixed positive-valuation cutoff,
only finitely many monomials remain, so this norm-preserving
contraction also applies to the completed series.

The support transitions just constructed carry a negative Laurent
cocycle to the same principal part, with its coefficients unchanged.
For a fixed all-negative $\alpha$, the integral coefficient bound
is $|c_\alpha|\leq b^{\sum_i(-\alpha_i)}$; the denominator
sublattice has the same bound multiplied by $|\varpi|$.
Along the cofinal system $b$ tends to one, so these bounds tend to
$1$ and $|\varpi|$.  Their colimit quotient is almost $R$.
We obtain the discrete
finite-sum principal part module tensored with $R$, shifted by $m$.

At a fixed module stage $r$, the perfected analytic scalars can
be computed by allowing a further scalar-root stage $h\geq r$.
This gives two indices: $r$ for the module transitions
\eqref{support:root-coefficient-transition}, and $h$ for the
fractional exponents in its scalar coefficients.  After the fixed
cutoff calculation every principal part is finite, so its exponents
have one common root denominator.  It appears at some finite $h$,
and increasing $r$ to at least $h$ places it on the diagonal.
That diagonal is cofinal.  The identity
$\theta_{r+1}\iota_r=F\theta_r$ identifies its algebraic Cech
transitions with the direct system on the left of
\eqref{support:principal-parts-comparison}.  The integral matrices
preserve both the coefficient lattices and their $\varpi$-multiple
submodules.

The comparison map itself is the map from algebraic Cech
fractions to the supported analytic Cech complex, permitting
arbitrarily small scalar losses before almostification.  It is
therefore independent of the chosen expression in monomials.
The residue change-of-variables formula proves coordinate
naturality.  On a compact continuous group parameter, uniform
polynomial approximation reduces the calculation modulo the
cutoff to locally constant finite principal parts.  Hence the
same map and contractions apply to the parameter complexes.
\end{proof}

Let $\mathbb D^\times$ be the puncture $x\ne0$ of the completed perfected
open labelled disc of $C_1$ in Proposition~\ref{prop:finite-rings}. At our
chosen sufficiently deep level
$U\subset P_t\times\ker(\GL_m(\Zp)\to\GL_m(\Fp))$, let
$\nu:\mathcal X/U\to\mathbb D^\times$ be the finite cover obtained by
analytically realizing the finite \'etale $C_1[1/x]$-algebra $B_U$.
The symbol $\langle\eta\rangle$ denotes
the retained one-dimensional volume representation; it is not
absorbed into a change of auxiliary action.

\begin{lemma}\label{support:positive-extension}
There is a natural almost equivalence on the ambient perfected
labelled disc
\begin{equation}\label{support:positive-extension-equivalence}
 \colim_r\mathscr M_r
 \simeq j_!\nu_*\mathscr A_{\mathcal X/U}
             \otimes\langle\eta\rangle,
 \qquad j:\mathbb D^\times\hookrightarrow\mathbb D^{m,\mathrm{perf}}.
\end{equation}
It respects finite-cover trace, changes of frame level, the
$P_n$-action, and the retained volume character.
\end{lemma}

\begin{proof}
Trivialize the finite volume frame, retaining its transformation
line.  By Lemma~\ref{lem:cartier-transfer}, the transpose is the
actual positive Frobenius on finite etale cover functions in this
frame.  In particular, the matrices here are not arbitrary
semilinear matrices.  Their root-stage bases are the roots of
these finite-cover functions.

First work on a bounded affinoid chart with $|x|$ bounded away
from zero.  A finite free basis norm and the intrinsic integral
norm on cover functions are equivalent, with a fixed two-sided
factor $C_0\geq1$.  At the $r$th root stage this factor is at most
$C_0^{1/Q^r}$.  The chosen basis is integral over $C_1$, so its
root basis is power bounded; thus there is an actual map from
these integral lattices into the intrinsic integral coefficient.
For every $\epsilon>0$, at sufficiently large $r$ the inverse
comparison loses norm at most $|\varpi|^{-\epsilon}$.  It follows
that the cokernel is almost zero.  Applying the same two-sided
estimates to the $\varpi$-multiples proves almost injectivity
after reduction.  Finally finite fractional polynomials are
dense in the completed perfection; modulo a positive-valuation
cutoff they are available after a common finite increase in
root stage.  This proves the desired equivalence on
$\mathbb D^\times$, including the completion comparison.

On the closed locus $x=0$, every transition of the free extended
lattices is zero, since its entries are divisible by the relevant
root of $x$.  Hence the direct limit has zero restriction to
that locus.  In terms of tubes, an element from a fixed stage
acquires a positive fractional power of $x$ at its next stage;
the latter basis is power bounded.  It therefore vanishes
modulo $\varpi$ on a sufficiently small tube.  The tube may
depend on the fixed stage.  This is enough for the direct-limit
restriction, and no common tube for all stages is needed.

The zero locus is generalizing closed, being the inverse limit
of its shrinking rational tubes.  Coefficient continuity and
\eqref{support:localization} now identify the direct limit with
extension by zero of the coefficient already computed on its
open complement.  This proves
\eqref{support:positive-extension-equivalence}.

The construction used the intrinsic finite-cover functions,
their actual Frobenius, and the residue/trace frame.  All these
maps commute with the finite-level operations of
Lemma~\ref{lem:cartier-transfer}.  Norm bounds serve only to
prove that these same maps are almost equivalences; they do
not choose a replacement action.  The estimates hold uniformly
on compact $P_n$-families by a common finite basis and pole
bound.  This proves the asserted naturalities.
\end{proof}

\subsection{The transpose comparison}

We use the ordinary constant-cochain action throughout this comparison.
On a $P_n$-action nerve, a locally constant function
$P_n^a\to\Fp$ gives a geometric cochain by multiplying the
coefficient unit $1$ on each clopen piece.  These maps commute with
the faces, degeneracies, and cup products.  Totalization thus gives
the structural action of $C^*_{\mathrm{cts}}(P_n;\Fp)$; on the
quotient presentation it is pullback from $BP_n$.  Supported fibers
inherit the action from the same nerve.  This specifies the action
before making any comparison of cohomology groups.

\begin{proposition}\label{prop:transpose-comparison}
For the stable lattice $P$ and sufficiently deep finite frame
level $U$ specified above, there is a natural equivalence
\begin{equation}\label{support:transpose-derived}
 \left(\colim_{S^\vee}
    \RGamma(P_n,P)^\vee\right)\otimes_kR
 \simeq
 \RGamma\bigl(U,\RGamma_c(Z,\mathscr A)
                     \otimes\langle\eta\rangle\bigr)[m+n^2].
\end{equation}
The dual of a compact cochain complex is its continuous
$k$-linear dual, with cochain degrees reversed.  In degree $-i$
this gives
\begin{equation}\label{support:transpose-cohomology}
 \left(\colim_{S^\vee}H^i(P_n,P)^\vee\right)\otimes_kR
 \simeq
 H^{m+n^2-i}\bigl(U,\RGamma_c(Z,\mathscr A)
                        \otimes\langle\eta\rangle\bigr).
\end{equation}
These are equivalences of modules for the action of the actual
constant $\Fp$-cochains of $P_n$, with the usual dual signs.
Transpose trace upon enlargement of a frame level corresponds
to pullback on the geometric side.
\end{proposition}

\begin{proof}
The comparison has four terms.  The two shifts come from group
duality and principal parts, respectively:
\begin{equation}\label{support:comparison-diagram}
\begin{tikzcd}[row sep=large]
 \left(\colim_{S^\vee}\RGamma(P_n,P)^\vee\right)\otimes_kR
   \arrow[d,"\text{group and residue duality}","\simeq"']\\
 \RGamma\!\left(P_n,
     (\colim_F H^m_{\mathfrak m}(M))\otimes_kR\right)[n^2]
   \arrow[d,"\text{principal parts and positive extension}","\simeq"']\\
 \RGamma\!\left(P_n,\RGamma_c(\mathcal X/U,\mathscr A)
                    \otimes\langle\eta\rangle\right)[m+n^2]
   \arrow[d,"\text{the common torsor nerve}","\simeq"']\\
 \RGamma\!\left(U,\RGamma_c(Z,\mathscr A)
                    \otimes\langle\eta\rangle\right)[m+n^2].
\end{tikzcd}
\end{equation}
We justify these arrows with their cochain actions.

The group $P_n$ has dimension $n^2$ and no $p$-torsion, since
$p-1>n$.  Its dualizing orientation is trivial: after a splitting
field, the adjoint action is conjugation on matrices and has
determinant one.  Continuous compact-group duality therefore gives
\[
 \RGamma(P_n,P)^\vee
 \simeq\RGamma(P_n,P^\vee)[n^2].
\]
One may compute this with a bounded finite projective resolution
over $k[[P_n]]$.  The compact and discrete models are compatible
by reduction to finite coefficients and inverse limit; the
transition maps on the terms are surjective.  This is also the
compact/discrete duality of Lemma~\ref{lem:compact:duality}.
The completed-module interpretation for the corresponding
continuous and solid group cochains is compatible with these
finite projective resolutions
\cite[Theorems~3.2 and~3.14(i), Proposition~3.21]{Tang}.

Apply \eqref{support:residue-dual} and
Lemma~\ref{support:principal-parts}.  The bounded resolution lets
us commute its finite sums and retracts with the filtered root
colimit.  Lemma~\ref{support:positive-extension} then identifies
the root system with the compact-support coefficient in the third
term of \eqref{support:comparison-diagram}.
There is no $P_n$-character on the constant volume frame; its
remaining character is the auxiliary one displayed here.

Finally use Theorem~\ref{thm:geometric-quotient}.  The two compact
torsor presentations of $[Z/U]$ have compatible cofinal compact
bounds.  At each fixed bound their common double nerve gives the
comparison; the bounded $P_n$ and $U$ resolutions in
Lemma~\ref{support:parameters} allow the support colimit to pass through
both group directions.  This is the last arrow of
\eqref{support:comparison-diagram}.
Lemma~\ref{support:parameters} gives this comparison with every
compact parameter and every support restriction needed in the
nerves.  This proves \eqref{support:transpose-derived} and its
cohomological form.

We check the extra module assertion.  On the original
presentation, the constant cochain algebra acts through the
structural map to $BP_n$.  Compact-group duality sends its cup
action to the cap/cup transpose with the canonical graded
signs.  The residue pairing, the Cech support maps, and the
finite-cover Frobenius comparison are all maps over that same
presentation.  They therefore commute with these operations.
The common double nerve carries the identical structural map
to $BP_n$, so the action transported to $[Z/U]$ is the actual
one, not an action chosen after comparing dimensions.  Finite
etale trace and its transpose are compatible with the same
residue pairing.  Enlarging a frame level hence gives geometric
pullback, as asserted.
\end{proof}

\begin{corollary}\label{support:finite-stable-image}
For every $i$, the stable transfer image
\[
 \bigcap_{r\geq0}S^rH^i(P_n,P)
\]
is finite dimensional over $k$.
\end{corollary}

\begin{proof}
Corollary~\ref{support:finite-invariants} and faithful scalar
extension applied to \eqref{support:transpose-cohomology} imply
that $\colim_{S^\vee}H^i(P_n,P)^\vee$ is finite dimensional.
Its compact dual is the inverse limit with transition $S$.
Projection from that inverse limit has image exactly the
intersection displayed above: compactness supplies compatible
preimages of every point in the intersection.  The intersection
is therefore a quotient of a finite-dimensional vector space.
\end{proof}
